% 每层的谱范数与层间激活 Lipschitz 因子；末层为线性输出。
\begin{tikzpicture}[foundation picture]
 \path[use as bounding box] (0,-4) rectangle (112,66);
 \node at (20,62) {$\bm x$};
 \node[foundation input,fill=FigureBlueFill,minimum width=32mm,minimum height=8mm] (wone) at (20,52)
  {$\bm W_1,\ s_1=\|\bm W_1\|_2$};
 \node[foundation model,fill=FigureRedFill,minimum width=32mm,minimum height=8mm] (wtwo) at (20,34)
  {$\bm W_2,\ s_2=\|\bm W_2\|_2$};
 \node[foundation output,fill=FigureYellowFill,minimum width=32mm,minimum height=8mm] (wlast) at (20,10)
  {$\bm W_L,\ s_L=\|\bm W_L\|_2$};
 \draw[foundation flow] (20,58)--(wone.north);
 \draw[foundation flow] (wone.south)--(wtwo.north);
 \node[anchor=east] at (16,43) {$\phi_1$};
 \node[anchor=west] at (24,43) {$a_1$};
 \draw[foundation flow] (wtwo.south)--(20,23);
 \node[anchor=east] at (16,25) {$\phi_2$};
 \node[anchor=west] at (24,25) {$a_2$};
 \node at (20,19) {$\vdots$};
 \draw[foundation flow] (20,16)--(wlast.north);
 \draw[foundation flow] (wlast.south)--(20,3.5);
 \node at (20,0) {$f_{\bm W}(\bm x)$};
 \node at (64,57) {线性层};
 \node at (86,57) {激活函数};
 \node (linear-factor) at (64,46) {$\displaystyle\prod_{\ell=1}^{L}s_\ell$};
 \node (activation-factor) at (86,46) {$\displaystyle\prod_{\ell=1}^{L-1}a_\ell$};
 \draw[FigureMuted,line width=1pt,decorate,decoration={brace,amplitude=3mm}]
  (41,56)--(41,0);
 \node (lip) at (80,28) {$\displaystyle\operatorname{Lip}(f_{\bm W})\leq
  \prod_{\ell=1}^{L}s_\ell\prod_{\ell=1}^{L-1}a_\ell$};
 \draw[foundation flow] (44.5,28)--(lip.west);
 \draw[foundation flow] (linear-factor.south)--(64,37)
  -| ($(lip.north west)!.60!(lip.north east)$);
 \draw[foundation flow] (activation-factor.south)--(86,37)
  -| ($(lip.north west)!.87!(lip.north east)$);
\end{tikzpicture}%
